\section{A cofactor-core bridge to the complete physical output}
\label{sec:tpc57-block-bridge}

The singleton floor becomes useful only after two support issues are kept
separate.  First, a cofactor block is a set of input coordinates, whereas
the target norm is the complete undeleted output.  A lower bound for the
grouped norm of the restricted block is not monotone under restoring the
omitted coordinates.  Second, whether a terminal-admissible coordinate is
physically active depends on the distinguished source coefficient.  This
section inserts both distinctions into the literal fixed-shift map and
identifies the exact weighted collision estimate that would close the
singleton route.

\subsection{The literal terminalized vector and a cofactor core}

Fix \(h_0\ne0\), one declared left source system, and the direct
residue/orbit slice inherited from TPC--53--56
\citep{WangTPC53,WangTPC55,WangTPC56}.  We use the exact terminal and
physical sets
\[
 \cU_{h_0}^{\rm phys}\subseteq
 \cU_{h_0}^{\rm term}\subseteq\cU^{\rm pre}
\]
and the complete output key
\begin{equation}
 q(u)=(\gamma,d_mr,j),
 \qquad u=(m,r,j,\gamma),
 \label{eq:tpc57-literal-output-key}
\end{equation}
from TPC--55.  Put
\begin{equation}
 x_u:=(D_{h_0}G_X^L)_u,
 \qquad u\in\cU_{h_0}^{\rm term}.
 \label{eq:tpc57-literal-terminalized-vector}
\end{equation}
By definition of \(\cU_{h_0}^{\rm phys}\), and because the terminal
multiplier \(\mu(r_u)\Lambda(p_ur_u)\) is nonzero on
\(\cU_{h_0}^{\rm term}\), one has
\begin{equation}
 x_u\ne0
 \quad\Longleftrightarrow\quad
 [G_X^L]_u\ne0
 \quad\Longleftrightarrow\quad
 u\in\cU_{h_0}^{\rm phys}.
 \label{eq:tpc57-terminal-physical-equivalence}
\end{equation}
Thus the active multiplicity of a physical output fiber is
\begin{equation}
 n_y^{\rm phys}
 :=\left|q^{-1}(y)\cap\cU_{h_0}^{\rm phys}\right|.
 \label{eq:tpc57-physical-fiber-degree}
\end{equation}
All norms below use raw counting measure.  In particular,
\begin{align}
 \cD_{h_0}(G_X^L)
 &=\sum_{u\in\cU_{h_0}^{\rm term}}|x_u|^2,
 \label{eq:tpc57-literal-terminal-diagonal}\\
 \cG_{h_0}(G_X^L)
 &=\sum_y\left|\sum_{u\in q^{-1}(y)}x_u\right|^2.
 \label{eq:tpc57-literal-complete-output}
\end{align}
The sum in \eqref{eq:tpc57-literal-complete-output} is the original
undeleted fixed-\(h_0\) synthesis, not a selected packet.

For a cofactor scale \(R>0\), define the pre-terminal and terminal core
sets
\begin{align}
 \cB_R^{\rm pre}
 &:=\{u\in\cU^{\rm pre}:R\le r_u\le2R\},
 \label{eq:tpc57-preterminal-cofactor-core}\\
 \cB_R^{\rm term}
 &:=\cB_R^{\rm pre}\cap\cU_{h_0}^{\rm term}.
 \label{eq:tpc57-terminal-cofactor-core}
\end{align}
Its terminal diagonal energy is
\begin{equation}
 \cD_R^{\rm core}
 :=\sum_{u\in\cB_R^{\rm term}}|x_u|^2.
 \label{eq:tpc57-core-terminal-energy}
\end{equation}
The rooted collision incidence of the core against the \emph{full}
physical vector is
\begin{equation}
 \cI_R^{\rm full}
 :=\sum_{u\in\cB_R^{\rm term}\cap\cU_{h_0}^{\rm phys}}
 |x_u|^2\bigl(n_{q(u)}^{\rm phys}-1\bigr).
 \label{eq:tpc57-core-full-physical-incidence}
\end{equation}
The multiplicity in \eqref{eq:tpc57-core-full-physical-incidence} is not
computed in the restricted block.  This is the essential point that makes
the resulting lower bound a statement about the complete output.

\begin{theorem}[Full-output cofactor-core singleton bridge]
\label{thm:tpc57-physical-core-bridge}
For the distinguished literal terminal vector,
\begin{equation}
 \boxed{
 \cG_{h_0}(G_X^L)
 \ge
 \bigl(\cD_R^{\rm core}-\cI_R^{\rm full}\bigr)_+.}
 \label{eq:tpc57-physical-core-floor}
\end{equation}
Consequently, if \(\cD_R^{\rm core}>0\) and
\begin{equation}
 \cI_R^{\rm full}
 \le(1-\delta_R)\cD_R^{\rm core}
 \qquad(0<\delta_R\le1),
 \label{eq:tpc57-core-incidence-target}
\end{equation}
then the complete undeleted output satisfies
\begin{equation}
 \boxed{
 \cG_{h_0}(G_X^L)
 \ge\delta_R\cD_R^{\rm core}.}
 \label{eq:tpc57-core-positive-output}
\end{equation}
\end{theorem}

\begin{proof}
Apply \cref{prop:tpc57-core-to-full-bridge} to the full terminalized vector
\eqref{eq:tpc57-literal-terminalized-vector} and the coordinate subset
\(\cB_R^{\rm term}\).  Equations
\eqref{eq:tpc57-terminal-physical-equivalence} and
\eqref{eq:tpc57-physical-fiber-degree} identify the abstract full active
multiplicity with \(n_y^{\rm phys}\), giving
\eqref{eq:tpc57-physical-core-floor}.  Substitution of
\eqref{eq:tpc57-core-incidence-target} gives
\eqref{eq:tpc57-core-positive-output}.
\end{proof}

No inequality between
\(\cG_{h_0}(G_X^L\one_{\cB_R^{\rm pre}})\) and
\(\cG_{h_0}(G_X^L)\) is used in the proof.  Coordinates outside the core
may cancel a restricted block sum.  The theorem avoids that invalid
monotonicity step by retaining only core coordinates that are singleton in
the full active fiber system.

\subsection{Every partner lies in one enlarged cofactor block}

Although \eqref{eq:tpc57-core-full-physical-incidence} uses full active
multiplicities, its possible partners are geometrically local.  Define
\begin{align}
 \cB_R^+
 &:=\{u\in\cU_{h_0}^{\rm term}:R/2\le r_u\le4R\},
 \label{eq:tpc57-enlarged-terminal-block}\\
 n_y^+(R)
 &:=\left|q^{-1}(y)\cap\cU_{h_0}^{\rm phys}\cap\cB_R^+\right|.
 \label{eq:tpc57-enlarged-physical-degree}
\end{align}

\begin{lemma}[Factor-exchange localization]
\label{lem:tpc57-factor-exchange-localization}
Let \(u\in\cB_R^{\rm term}\), and suppose
\(v\in\cU_{h_0}^{\rm term}\) satisfies \(q(v)=q(u)\).  Then
\begin{equation}
 R/2\le r_v\le4R.
 \label{eq:tpc57-partner-cofactor-range}
\end{equation}
In particular, if a physical fiber meets the core, then its complete
physical multiplicity is already visible in the enlarged block:
\begin{equation}
 n_{q(u)}^{\rm phys}=n_{q(u)}^+(R)
 \qquad
 \bigl(u\in\cB_R^{\rm term}\cap\cU_{h_0}^{\rm phys}\bigr).
 \label{eq:tpc57-full-degree-enlarged-identity}
\end{equation}
\end{lemma}

\begin{proof}
Write
\[
 u=(d\ell,r,j,\gamma),
 \qquad
 v=(d'\ell',r',j,\gamma).
\]
Equality of the complete output keys gives the exact factor exchange
\begin{equation}
 dr=d'r'.
 \label{eq:tpc57-factor-exchange-equation}
\end{equation}
Both source squarefree factors lie in the inherited dyadic interval
\([D,2D]\).  Hence
\[
 r'=\frac d{d'}r,
 \qquad
 \frac12\le\frac d{d'}\le2.
\]
Together with \(R\le r\le2R\), this proves
\eqref{eq:tpc57-partner-cofactor-range}.  Every full physical partner of a
physical core atom is therefore counted by
\eqref{eq:tpc57-enlarged-physical-degree}, proving
\eqref{eq:tpc57-full-degree-enlarged-identity}.
\end{proof}

The constants \(1/2\) and \(4\) come only from the declared dyadic
squarefree support.  Thus enlarging the core creates an \(O(1)\) cover and
has no positive polynomial cost.  More explicitly,
\begin{equation}
 \boxed{
 \cI_R^{\rm full}
 =\sum_{\substack{u\in\cU_{h_0}^{\rm phys}\cap\cB_R^{\rm term}\\
                  v\in\cU_{h_0}^{\rm phys}\cap\cB_R^+\\
                  v\ne u,\ q(v)=q(u)}}|x_u|^2.}
 \label{eq:tpc57-enlarged-rooted-pair-identity}
\end{equation}
This is an ordered pair sum with the terminal energy placed on the core
root \(u\).  It is not an unweighted count of factor exchanges.

\subsection{Row-factorized form of the exact missing estimate}

The exact literal interface
\eqref{eq:tpc57-exact-row-factorization} factors every terminal amplitude as
\begin{equation}
 x_u=\Omega_{q(u)}\kappa_u\zeta_{m_u}.
 \label{eq:tpc57-literal-row-factorization}
\end{equation}
Here \(\kappa_u\) is the product of the shaped terminal logarithm, direct
source shapes, mask, and every other displayed factor in the remaining
amplitude except \(\zeta_{m_u}\).  Put
\begin{equation}
 c_u:=\Omega_{q(u)}\kappa_u
 \label{eq:tpc57-baseline-physical-amplitude}
\end{equation}
and define the baseline-nondegenerate carrier
\begin{equation}
 \cU_{h_0}^{\circ}
 :=\{u\in\cU_{h_0}^{\rm term}:c_u\ne0\}.
 \label{eq:tpc57-baseline-terminal-carrier}
\end{equation}
Then
\begin{equation}
 x_u=c_u\zeta_{m_u},
 \qquad
 u\in\cU_{h_0}^{\circ},
 \label{eq:tpc57-baseline-times-row-coefficient}
\end{equation}
and a baseline edge is physically active precisely when its row coefficient
is nonzero.  Distinct terminal atoms in a fixed output fiber have distinct
source rows by the prime-resolved injection of TPC--55.  Thus a fixed pair
\((m,y)\) supports at most one baseline edge, as required by
\cref{thm:tpc57-row-minimax}.

For each row define its core baseline energy
\begin{equation}
 W_{m,R}
 :=\sum_{\substack{u\in\cU_{h_0}^{\circ}\cap\cB_R^{\rm term}\\
                   m_u=m}}|c_u|^2.
 \label{eq:tpc57-core-row-weight}
\end{equation}
On the enlarged baseline carrier put
\begin{equation}
 d_y^+(R)
 :=\left|q^{-1}(y)\cap\cU_{h_0}^{\circ}\cap\cB_R^+\right|,
 \label{eq:tpc57-enlarged-baseline-degree}
\end{equation}
and define the support-only rooted row incidence
\begin{equation}
 J_{m,1}^{R,+}
 :=\sum_{\substack{u\in\cU_{h_0}^{\circ}\cap\cB_R^{\rm term}\\
                   m_u=m}}
 |c_u|^2\bigl(d_{q(u)}^+(R)-1\bigr).
 \label{eq:tpc57-enlarged-baseline-row-incidence}
\end{equation}
For the distinguished coefficient vector, let
\begin{equation}
 n_y^+(R;\zeta)
 :=\left|\{v\in q^{-1}(y)\cap\cU_{h_0}^{\circ}\cap\cB_R^+:
                   \zeta_{m_v}\ne0\}\right|
 \label{eq:tpc57-enlarged-active-row-degree}
\end{equation}
and
\begin{equation}
 \mathfrak J_{m,1}^{R,+}(\zeta)
 :=\one_{\{\zeta_m\ne0\}}
   \sum_{\substack{u\in\cU_{h_0}^{\circ}\cap\cB_R^{\rm term}\\
                   m_u=m}}
 |c_u|^2\bigl(n_{q(u)}^+(R;\zeta)-1\bigr).
 \label{eq:tpc57-enlarged-physical-row-incidence}
\end{equation}
The row indicator makes the active-root restriction literal.  In
particular, an inactive row contributes zero rather than a spurious negative
value from \(n_{q(u)}^+(R;\zeta)-1\).

\begin{proposition}[Exact distinguished incidence and its baseline majorant]
\label{prop:tpc57-weighted-incidence-factorization}
The core diagonal and full rooted collision incidence satisfy
\begin{align}
 \cD_R^{\rm core}
 &=\sum_m|\zeta_m|^2W_{m,R},
 \label{eq:tpc57-weighted-core-energy}\\
 \cI_R^{\rm full}
 &=\sum_m|\zeta_m|^2
   \mathfrak J_{m,1}^{R,+}(\zeta),
 \label{eq:tpc57-exact-weighted-physical-incidence}\\
 0\le\mathfrak J_{m,1}^{R,+}(\zeta)
 &\le J_{m,1}^{R,+}.
 \label{eq:tpc57-physical-baseline-incidence-order}
\end{align}
Consequently the coefficient-weighted support estimate
\begin{equation}
 \sum_m|\zeta_m|^2J_{m,1}^{R,+}
 \le(1-\delta_R)\sum_m|\zeta_m|^2W_{m,R}
 \label{eq:tpc57-baseline-weighted-collision-target}
\end{equation}
is sufficient for the undeleted lower bound
\eqref{eq:tpc57-core-positive-output}.
\end{proposition}

\begin{proof}
Equation \eqref{eq:tpc57-weighted-core-energy} follows by substituting
\eqref{eq:tpc57-baseline-times-row-coefficient} into
\eqref{eq:tpc57-core-terminal-energy}.  For an active core root \(u\),
\cref{lem:tpc57-factor-exchange-localization} identifies its full active
degree with \(n_{q(u)}^+(R;\zeta)\).  Substitution into
\eqref{eq:tpc57-core-full-physical-incidence} and grouping the roots by
their source row proves \eqref{eq:tpc57-exact-weighted-physical-incidence}.
Since active enlarged edges form a subset of the baseline enlarged carrier,
\(n_y^+(R;\zeta)\le d_y^+(R)\), proving
\eqref{eq:tpc57-physical-baseline-incidence-order}.  Therefore
\eqref{eq:tpc57-baseline-weighted-collision-target} implies
\eqref{eq:tpc57-core-incidence-target}, and
\cref{thm:tpc57-physical-core-bridge} finishes the proof.
\end{proof}

The exact weakest first-moment target is
\eqref{eq:tpc57-exact-weighted-physical-incidence} with
\(\mathfrak J_{m,1}^{R,+}(\zeta)\).  Replacing it by the baseline quantity
\(J_{m,1}^{R,+}\) is a stronger coefficient-independent sufficient
condition.  Neither target follows from a maximal-degree upper bound: a
carrier may consist entirely of doubleton fibers.

\subsection{Activation and the prospective loss ledger}

To state the route without hiding an occupancy assertion, let
\(\cE_{\rm par}\) be the parent physical atomic energy before selection of
the declared literal packet.  Define the selected pre-terminal energy and
the core pre-terminal energy by
\begin{align}
 \cE_{\rm pre}(G_X^L)
 &:=\sum_{u\in\cU^{\rm pre}}|[G_X^L]_u|^2,
 \label{eq:tpc57-selected-preterminal-energy}\\
 \cE_R^{\rm pre}(G_X^L)
 &:=\sum_{u\in\cB_R^{\rm pre}}|[G_X^L]_u|^2.
 \label{eq:tpc57-core-preterminal-energy}
\end{align}
Suppose, along a sequence of scales,
\begin{align}
 \cE_{\rm pre}(G_X^L)
 &\ge X^{-\lambda_{\rm sel}+o(1)}\cE_{\rm par},
 \label{eq:tpc57-parent-selection-hypothesis}\\
 \cE_R^{\rm pre}(G_X^L)
 &\ge X^{-\lambda_{\rm blk}+o(1)}
       \cE_{\rm pre}(G_X^L),
 \label{eq:tpc57-core-block-occupancy-hypothesis}\\
 \cD_R^{\rm core}
 &\ge X^{-\lambda_{\rm term}+o(1)}
       \cE_R^{\rm pre}(G_X^L),
 \label{eq:tpc57-core-terminal-activation-hypothesis}\\
 \delta_R
 &\ge X^{-\lambda_{\rm sing}+o(1)},
 \qquad
 \cI_R^{\rm full}\le(1-\delta_R)\cD_R^{\rm core}.
 \label{eq:tpc57-core-singleton-fraction-hypothesis}
\end{align}
Then the exact bridge gives
\begin{equation}
 \boxed{
 \cG_{h_0}(G_X^L)
 \ge
 X^{-(\lambda_{\rm sel}+\lambda_{\rm blk}
       +\lambda_{\rm term}+\lambda_{\rm sing})+o(1)}
 \cE_{\rm par}.}
 \label{eq:tpc57-prospective-activation-ledger}
\end{equation}

Equation \eqref{eq:tpc57-prospective-activation-ledger} is a conditional
composition, not an assertion that any of
\eqref{eq:tpc57-parent-selection-hypothesis}--
\eqref{eq:tpc57-core-singleton-fraction-hypothesis} holds for the inherited
packet.  The new singleton gate is exactly
\(\lambda_{\rm sing}\).  There is no additional coherence exponent after
\eqref{eq:tpc57-core-singleton-fraction-hypothesis}, because the surviving
core energy lies in singleton fibers of the complete output.  Native-Gram,
boundary, and tail comparisons remain separate downstream operations.
